論理を辿って一歩づつ考える能力は、人類の知的活動の源泉です。
LLMは膨大な知識に基づき「それっぽいこと」を言うことはできますが、
「論理的に正しい」推論を行うことができるでしょうか？


我々は、LLMに対して論理推論を教えるための研究を進めてきています。

この目的のためにはまず、「論理的に正しいとはどういうことか？」という問いに立ち返る必要があります．
そこで我々は、論理的な推論を探求する学問である記号論理学の知見を借りつつ、論理推論の特徴を明確にしていきます。
論理推論は、「"A"という事実と、"AならばB"という事実があった時に、"B"という事実を導いてよい」というような、「推論の規則」を用いて、新しい事実を導きます。
このような推論規則は多様な種類が考えられますが、その中でも「基礎的」なものは、少数に限られることが知られています。
また、この規則の"A"や"B"としては「任意の事実」が許される、という点にも特徴があります。

次に、LLMに論理推論を教えるための手立てを考える必要があります。
LLMはデータから物事を学習する機械なので、論理推論を学ぶためのデータが必要となります。
我々はこのためのデータとして、「論理推論問題」を作ることにしました。
与えられた事実を基に論理を辿ることで、最終的な回答に辿り着きます。
そして、論理推論の特徴をきちんと捉えるための、問題の設計指針を定めました。
例えば「基礎的な推論規則を組合わせることで解けるようにする」「"A"や"B"に入る事実が任意であることを教えるために、様々な事実を用いて問題を構成する」などです。

次に、設計指針を反映した問題を大量に用意する必要があります。
このために我々は、この問題を自動生成するアルゴリズムを開発しました。
推論規則を多段階積み重ねることで問題の骨格を作り、語彙・自然言語表現テンプレートを用いて作成した事実によって肉付けします。

この問題をLLMに大量に解かせたところ、論理推論タスクでの性能が大幅に向上しました。
興味深いこととして、数学や科学・コーディング等、幅広い推論タスクでの性能向上も得られました。
この理由は、論理推論というのは、他の推論の基礎となるものだからだと考えます。

本研究は、2022年頃、LLMが登場する少し前くらいから始めて、現在まで継続しています。
我々は、学習データの他にも、LLMが論理的に正しい推論を行えるかどうかを評価するベンチマークも作っています。
単に「それっぽいことを言う」だけでは正解できなような問題設定・評価設計を作りました。